\chapter{Εισαγωγή}
\pagenumbering{arabic}%\pagestyle{mythesis}
\label{ch:intro}
\hrule\vspace{0.2in}
\begin{itemize}
  \item[\ref{sec:11}] Στόχοι
  \item[\ref{sec:12}] Δομή της Διατριβής
\end{itemize}
\hrule \vspace{0.2in}



\section{Στόχοι}\label{sec:11}
Εδώ αρχίζει η κανονική αρίθμηση των σελίδων. Παραπάνω η αρίθμηση
είναι i, ii, iii, … όπως φαίνεται.

Συνήθως αυτό το κεφάλαιο (ΕΙΣΑΓΩΓΗ) πρέπει να υπάρχει. Στο
περιεχόμενό του πρέπει να συμπεριλαμβάνονται και τα ακόλουθα:

\y Στόχοι. Για παράδειγμα: Ο συνολικός στόχος της διατριβής είναι
… Οι επί μέρους στόχοι :
\begin{itemize}
\item Να αναλυθεί η συμπεριφορά ...
\item Η ανάπτυξη  ...
\item Να κατηγοριοποιηθούν  ...
\item Να επιβεβαιωθούν  ...
\end{itemize}




\section{Δομή της Διατριβής}\label{sec:12}
Η διατριβή περιέχει ν κεφάλαια: Το Κεφάλαιο 1 ….


\newpage
Το Σχήμα~\ref{fig1} δείχνει το Πρώτο Σχήμα του Πρώτου Κεφαλαίου.

\begin{figure}[t]
 \centerline{ {\includegraphics[scale=0.9]{bird.eps}}
              \hskip 2.3in
              {\includegraphics[scale=0.9]{bird.eps}}}
 \centerline{(a) \hskip 2.5in (b) \hskip 0.1in}
 \centering
 \vspace{0.0cm}
 \caption{(a)~Ένα Σχήμα\  και \  (b)~δίπλα ένα άλλο Σχήμα.}\label{fig1}
\end{figure}
